% !TEX root = ../monograph.tex
% May be placed in Chapter13 after the graded quotient estimate.
% This fragment uses neither phase decomposition nor stress-tensor Ward identities.
\section{Local replacement of a cluster on the lattice}
\label{ch18:lattice-products}

We prove a local replacement estimate before using any phase decomposition.
It will also apply to the undecomposed correlations, before any decomposition into phases, with a fixed exterior word.
Observable source variables and local polynomial extraction are standard
renormalization-group tools; see, for example,
\cite{BauerschmidtBrydgesSlade2015}. The quantitative estimates below concern
the block kernels, activity norms and quotient fields constructed
here; that reference does not give these particular estimates.

For $T\to\infty$, write $f\preccurlyeq e^{cT}$ if, for every
$\varepsilon>0$, eventually $|f|\le e^{(c+\varepsilon)T}$ on an ordered
cutoff and volume tail, uniformly on the specified bounded label sets
and geometries. A factor $e^{o(T)}$ is called subpower in this paragraph.
Choose scales $l=l(T)$ with $\tau_l/T\to\theta$, where
$\tau_l=\log(1/a_l)$. The sequence of scales has $\log L_l=o(\tau_l)$, so
first-crossing scales give any positive $\theta$. For $\theta=0$ one may
first cross $\sqrt T$.

\subsection{Signed local transfer}

A local degree-$t$ density with insertions consists of a column coefficient
$x_l$ and a remainder of norm $r_l$ whose extracted coefficient vanishes. Its dimensionless
normalization obeys
$H_l=L_l^t\mathcal B_l H_{l+1}$; the dimensionful density with insertions has
no additional factor. These conventions are those of
Section~\ref{ch11:one-graph}. Keep the full prescribed large-field set and
its normalization by the square root of the large-field bound in this remainder norm.

The backward comparison of Section~\ref{ch11:trial-descent} applies
on any sufficiently late interval $[k_*,m]$, not only from the original
cutoff scale. In detail, a remainder of size $q$ with vanishing extracted coefficient has
output coefficient and remainder at most $L_l^{-2/5}q$.
The regular term uses its zero radial Taylor polynomial through $t$;
the large-field term retains half the suppression of its original
prescribed large-field set. Subtracting its output column restores the vanishing extracted coefficient and
adds at most $(1+2C_l)L_l^{-2/5}q$. The lower bound on the block ratio makes this at
most $L_l^{-3/8}q$. This recentering is an equality of
densities with insertions, including the terms with prescribed large-field sets, and uses linearity rather
than division by an activity.

Add a free column parameter at the finest scale and choose it so that
the transferred residual coefficient vanishes at $k_*$. Invertibility
of the column map permits this on the same history. The cone calculation
in Section~\ref{ch11:trial-descent} then gives
\begin{equation}
 r_l+|x_l|_1\le4(1+C_m)q_{\rm fin}(a_m/a_l)^{1/8},
 \qquad r_m\le2(1+C_m)q_{\rm fin}.
 \label{ch18:signed-transport}
\end{equation}
Here $q_{\rm fin}$ is the combined original coefficient and remainder
bound, before the affine column subtraction. To recall the numerical
point, with $Y=|\mathsf d_lx_{l+1}|_1$, $U=|x_l|_1$ and
$\epsilon_l=L_l^{-3/8}$ the two inequalities are
$Y\le K_l[U+\epsilon_l(Y+r_{l+1})]$ and
$r_l\le K_lY+\epsilon_lr_{l+1}$.
The cone $32(1+C_l)^2U\le r_l$ gives
$Y\le4K_l\epsilon_lr_{l+1}$ and $L_l^{1/8}r_l\le r_{l+1}$.
It starts with $x_{k_*}=0$. At the finest scale absorb the subtracted
column bound to obtain the second inequality, and iterate for the first.
The sign of the local difference has been retained throughout.

\subsection{Grouping insertion points and retaining their supports}

Use the square-zero source algebra of Chapter~\ref{ch11:insertions}.
The total number of copies may be arbitrary, but a reflected cell
contains a fixed number $n$ of insertion points. A group is a set of
these insertion points with a common representative. Grouping is bookkeeping for
supports and source weights: original variables and the coefficient of every source monomial
remain in the expansion. In particular partial modes of a group are
not discarded, and grouping is not asserted to be an algebra
homomorphism replacing their product by a new variable.

Recall that $c_B(X)$ counts the radius-$B$ balls in a smallest connected covering of $X$, and $\sigma$ is its exponential support weight.
If the support $X$ of a term with insertions has to contain new representatives joined
by at most $h$ covering balls per insertion point, then
\[
 c_B(X\cup\{y_i:i\in I\})\le c_B(X)+h|I|.
\]
Shrinking each affected weight by $\eta e^{-\sigma h}$ absorbs this
support enlargement. The incidence norm contributes a second factor: a new
incidence can be attributed to an old insertion point on its joining path.
If at most $\mathfrak N$ representatives occupy a covering ball, there
are at most $(1+\mathfrak N)n$ assignments per representative.
Take $\eta\le[(1+\mathfrak N)n]^{-1}$. Supports without insertions are
unchanged. This is a termwise estimate on the full
series with insertions, hence is uniform in its number of replica cells.

At each scale retain a maximal $B_l^2$-separated set of representatives,
keeping the representative of the center, and join each discarded
representative to a retained neighbor. A merger decreases the number
of groups; there are at most $n-1$ mergers within a cell.
A deliberately generous joining bound is
\[
 h_l=16[1+L_l(B_l+B_{l+1}+1)^2]^4.
\]
Impose $(\sigma_{l+1}+1)h_l=o(\tau_l)$ by choosing the lower bounds on the reference coupling at each stage.
At fixed stage $h_l$ is a polynomial in $\log\widehat\beta_l$
with fixed ratio, so this is a finite admissible request.
The occupancy factors have the same property. The minimum
$\alpha_j\le1$ of all weight-reduction factors for support enlargement on $m\le l<j$
therefore has $-\log\alpha_j=o(\tau_j)$: split the maximum
cost into a fixed initial segment and a tail bounded by
$\varepsilon\tau_j$, and then let $\varepsilon\downarrow0$.

Between mergers a single group loses at most the factor $6L_l$ in
its weight. Here the incoming regular contribution is carried
\emph{without} Taylor subtraction. Its full factor for the change of cover radius
is $T_l\le L_l$, as in \eqref{ch11:cover-transfer}; it
must not be assigned the contraction factor for a vanishing extracted coefficient.
If its incoming mass is $\mathfrak m$, the other local terms with vanishing extracted coefficient have total mass $p_l$ with $2p_l\le\mathfrak m$
after increasing the lower bound on the reference coupling, and the large-field terms have normalized mass
$f_l\le\mathfrak m$. These are the termwise local bounds used
above, retaining the original events. Resolving real and imaginary
parts contributes at most a factor two, so
\[
 2(p_l+T_l\mathfrak m)+f_l
             \le(2L_l+2)\mathfrak m\le6L_l\mathfrak m
             \le L_l^4\mathfrak m.
\]
The last inequality holds for $L_l\ge3$. It applies in the
submultiplicative norm in the source variables to a whole incoming joint factor;
it does not multiply separate bounds for correlated insertions.
With $g$ groups left, the product of their remaining source weights is consequently
at least
\begin{equation}
 (a_j/a_l)^{4n}(\alpha_j/8)^{n(n-g)}\prod_iw_i.
 \label{ch18:group-weight}
\end{equation}
Each step loses at most the factor $L_l^{-4n}$, and a merger uses at most
$n$ of the weight-reduction factors for support enlargement. This proves the formula by
induction, including all data on supports and prescribed large-field sets.

\subsection{Extraction with a freely chosen residual weight}

Put
\[
 \delta_l^{\rm mark}=\widehat\beta_{l+1}^{-2},\quad
 \phi_l=\widehat\beta_{l+1}^{-1},\quad
 \Xi_l=e^{\sqrt{\widehat\beta_{l+1}}},\quad
 y_{t,l}=\widehat\beta_l^{-\nu_t}.
\]
Choose the fixed integer $\nu_t$ to absorb the finite polynomial and
attachment losses of the insertions at degree $t$. The stage lower bounds on the reference coupling ensure
$y_{t,l}C_l\le1$ and
$\phi_l L_l^{t+1}y_{t,l}\le y_{t,l+1}$.
These are finite requests on the lower bound on the reference coupling after the ratio is fixed.
All their logarithms, as well as $\log\Xi_l$, are $o(\tau_l)$.

For a fixed exterior reflection cell at scale $k$, let
\[
 k_* =\min\{l\ge k:B_l^2\le a_k/a_l\},\qquad
 N_k=\exp\{(a_k/a_{k_*})^4q_{k_*}^{\rm site}+5\}.
\]
The stopping-scale argument of
\eqref{ch13:sharp-schedule} gives $\tau_{k_*}/\tau_k\to1$
and $\log N_k=o(\tau_{k_*})$. Indeed if $k_*=l+1>k$,
minimality gives $a_k/a_{k_*}\le L_lB_l^2$; the last condition
of that equation bounds both the normalization exponent and
$\tau_{k_*}-\tau_k$ by $o(\tau_{k_*})$.
When $\tau_k=o(T)$ the same proof gives $\log N_k=o(T)$.

Suppose the cluster has become one group at scale $m+1$, with
product $\omega$ of its remaining source weights. Set
\[
 q_m^{\rm ext}=4\Xi_ma_{m+1}^t/\omega,\qquad
                    E_m=8(1+C_m)q_m^{\rm ext}.
\]
Taylor extraction through radial degree $t$ gives coefficient and
complement each bounded by $q_m^{\rm ext}$. To verify this, divide
by this product of weights, apply the integral Taylor formula in the
local derivative norm, and sum the finite coefficient conversions
and the polynomial factors in the support size with their positive margin in the exponent. Their
bounds are fixed stage constants and powers of logarithmic radii
and inverse thresholds, hence eventually at most $\Xi_m$.
The factor $a_{m+1}^t$ converts the dimensionful density with insertions to the
normalized degree-$t$ density. Large-field terms keep their
original normalized majorants. Extraction of the coefficient of
all insertions in a group is a linear operation on the exact density;
products of partial modes are retained before taking it. Each
term is assigned the union of its supports and original prescribed large-field set,
so this operation never divides by the exact-cover gas of the density without insertions.

Choose a residual weight $\varpi>0$ at $k_*$ and thereafter
$\varpi_{l+1}=\phi_l\varpi_l$. The two sufficient inequalities are
\begin{align}
 \varpi(a_m/a_{k_*})^{t+1}y_{t,m}/y_{t,k_*}
        &\le\delta_m^{\rm mark}\omega/(16\Xi_m),\notag\\
 32(1+C_m)\varpi a_{k_*}^{-t}E_m(a_m/a_{k_*})^{1/8}
        &\le y_{t,k_*}/2.
 \label{ch18:residual-weights}
\end{align}
The first absorbs the local polynomial coefficient bound and the second the bound for its signed transfer.
Indeed the recursion for these bounds gives
$\varpi_l\le\varpi(a_l/a_{k_*})^{t+1}y_{t,l}/y_{t,k_*}$.
Combining this with \eqref{ch18:signed-transport} and the dimensionful
factor $a_l^{-t}$, the ratio of required to allowed weight at
scale $l$ is at most its value at $k_*$ times
$(a_l/a_{k_*})^{7/8}\le1$. Thus no intermediate scale has been
left unchecked.

The extracted matched field has coefficient norm at scale $m$
at most $2E_m$. The interior residual has chessboard seminorm
at most $2N_k/\varpi$. To see the latter on arbitrarily many
copies, put one residual in every reflected cell before integrating.
Taking the coefficient of the required source monomial in the sum with insertions contributes the factor $\varpi^{-1}$ per
cell; the normalized integral bound \eqref{ch11:root-bound}
contributes the factor $N_k$ per cell. Apply this separately to the original local
difference and its compensating column, and use the chessboard
triangle inequality. The rooted component estimates and bounds
used in this calculation depend on the number of insertions per cell, not on the
number of cells. All the terms involving at least two insertions generated during transfer
are included in these sums. The estimate is therefore stronger
than a bound on one value of a correlation functional.

\begin{proposition}[Lattice cluster replacement]
\label{ch18:lattice-replacement}
Let a cluster contain $n$ fields of total degree $W$. Choose
$0\le\theta_-<\theta<\theta_+$ and scales $k,m,j$ with
$\tau_k/T\to\theta_-$, $\tau_m/T\to\theta$, and $\tau_j/T\to\theta_+$. The cluster lies inside the centered block of scale $m+1$,
its points lie on the scale-$j$ grid with separation at least
$B_j^2a_j$, and its dependence sets lie inside the scale-$k$
cell. Put $Z=W\theta_++4n(\theta_+-\theta)$.
For each fixed auxiliary degree $t$ there are output labels
$v_{N,p}^{\rm out}\in V_t$, on the actual histories, satisfying
\begin{align}
 |x_m(v_{N,p}^{\rm out})|_1
       &\preccurlyeq e^{(Z-t\theta)T},\notag\\
 \left\|\prod_iO_{N,p}(v_i,x_i)
                  -O_{N,p}(t,v_{N,p}^{\rm out};0)\right\|_{\mathcal C_k}
       &\preccurlyeq e^{[Z-(t+1/16)(\theta-\theta_-)]T}.
 \label{ch18:lattice-estimate}
\end{align}
The difference is interior. Constants and thresholds are uniform in
cutoff and period on the ordered tails and in the number of copies
of this fixed local cluster.
\end{proposition}
\begin{proof}
The weights of normalized local insertions in
\eqref{ch13:weight-size} give $\prod_iw_i=e^{-W\theta_+T+o(T)}$.
Keep the coefficients of all incoming source monomials above $j$; separation permits their
bounds as terms with insertions but does not assert statistical independence. The grouping
estimate gives $\omega=e^{-ZT+o(T)}$ at $m+1$, and therefore
$E_m=e^{(Z-t\theta)T+o(T)}$.
Set
\[
 \varpi=\exp\{[(t+1/16)(\theta-\theta_-)-Z]T\}.
\]
In the first inequality of \eqref{ch18:residual-weights}, the
left exponent is $-Z-15(\theta-\theta_-)/16$, strictly smaller
than the right exponent $-Z$. In the second, the full left
exponent is $-(\theta-\theta_-)/16<0$ while the right side
has exponent zero. Both strict margins absorb every subpower
factor. The extraction estimate and $\log N_k=o(T)$ give the
claim. All support shifts take place within the specified cell;
reflections and conjugations retain the same bounds.
\end{proof}

\subsection{A consequence for undecomposed correlations without phases}
\label{ch18:mixed-replacement}

Fix finite correlation reconstruction maps onto the filtered spaces $F_D$.
More explicitly, choose a basis in each successive quotient and
extend it to $F_D$, with dual coordinates realized by finite
linear combinations of separated probe correlations. Arrange that
a degree-$i$ coordinate vanishes on $F_{i-1}$, and that reconstruction
fixes every field in $F_D$. Keep each old coordinate functional as
its original correlation formula on every larger field space. When
lifting a new quotient basis, subtract its values under all the old
coordinates, so those coordinates annihilate the new basis vectors.
This is possible by finite-dimensional separation by all spectator
correlations and successive subtraction of lower degrees. Consequently
the degree-$D$ reconstruction of a field in $F_t$, $t>D$, is exactly its
low-degree component in this fixed biorthogonal system. Denote the resulting reconstruction of a word by
$\mathcal P_Dw$. It depends on the cluster, not on its later spectator.
No rotation average or phase measure is used here.

\begin{proposition}[Local replacement for undecomposed correlations]
\label{ch18:mixed-bound}
Let $x_i=o+r z_i$, with the $z_i$ in a fixed compact separated
configuration set, and let $W$ be the total cluster degree.
For any fixed separated exterior word $v$, any $D$ and $\varepsilon>0$,
\begin{align}
 |b_\alpha(w;o)|&\le C r^{|\alpha|-W-\varepsilon},\notag\\
 |S(wv)-S((\mathcal P_D(\tau_{-o}w))(o)v)|
                   &\le C r^{D+1-W-\varepsilon}.
 \label{ch18:mixed-product-rate}
\end{align}
Here $b_\alpha$, $|\alpha|\le D$, are the fixed probe coordinates. The same field
$\mathcal P_D(\tau_{-o}w)$ works for every exterior word. The
estimates are uniform on fixed bounded label sets and separated
geometries, and hold with any fixed finite derivatives in the positions,
with their increased scaling dimensions.
\end{proposition}
\begin{proof}
First use generic points of the scale-$j$ grid. Put $T=\log(1/r)$ and
choose $0<\theta_-<\theta<1<\theta_+$. The cluster then fits
the block and separation conditions for large $T$.
For each fixed geometry, the output labels in
\eqref{ch18:lattice-estimate} are bounded at fixed $m$ and hence
at the reference scale by the inverse-column bound. Take a
subsequence along which the coefficients converge, and call its field $A_t$.
The chessboard inequality bounds its residual paired with any
fixed exterior word by exponent
$Z-(t+1/16)(\theta-\theta_-)$. The same subsequence works for all
spectators, since the coefficient limit was taken before pairing.

The graded quotient estimate \eqref{ch18:graded-distance} shows
that a degree-$i$ coordinate of $A_t$ has exponent $Z-i\theta$.
Its coordinate functional vanishes on $F_{i-1}$; this is where
an ungraded coefficient estimate would lose the desired power.
Choose $t$ so large that
$(t+1/16)(\theta-\theta_-)\ge(D+1)\theta$.
This also bounds every raw coordinate being evaluated, whose degree is
at most $D$. We do not require the residual to bound all coordinates
through the auxiliary degree $t$. Pairing the residual
with the fixed probes bounds the coordinate errors. Reconstruct
$A_t$ and omit degrees $i>D$. Each omitted coordinate has exponent
at most $Z-(D+1)\theta$, and its fixed exterior correlations are
bounded. This gives that exponent for the raw replacement error,
and $Z-i\theta$ for each raw coordinate.

For completeness the grid restriction can be removed without
assuming differentiability of the selected coefficients of $A_t$.
Use the fixed correlation reconstruction instead. Its coefficients and
the pairings of its residual are smooth functions of the cluster positions.
Their derivatives replace an input by a field of degree one higher.
The sharp separated bounds, using fixed cells for exterior probes
and shrinking cells only for cluster factors, give a crude exponent
$W+1$: choose total loss one in the exponent. Thus one can take $c_0=1$.
The grid exponent is $\theta_+W+c_1$, where
$c_1=4n(\theta_+-\theta)-(D+1)\theta$ for the residual, and
$c_1=4n(\theta_+-\theta)-i\theta$ for a degree-$i$ coordinate.
These numbers are independent of $W$ with $n,D,i$ fixed.
Multiplicative constants may depend on each of the finitely many
higher derivative orders used in the following induction. For all sufficiently large total
degrees the crude estimate implies the grid target, because
$\theta_+>1$. Use downward induction on the total degree. Moving to a generic grid
configuration contributes the factor $O(a_j)=e^{-\theta_+T+o(T)}$; the derivative
bound at degree $W+1$ gives precisely the target exponent at
$W$. The connecting paths remain separated. Only finitely many
higher derivatives are needed for each starting degree, so their
constants and thresholds can be chosen together. Finally let
$\theta\uparrow1$ and $\theta_+\downarrow1$ and absorb the
remaining loss into the prescribed $\varepsilon$.
This proves the estimate for undecomposed correlations using neither phase symmetry nor
an operator-product theorem.
\end{proof}
